\begin{table*}[!tbp]
\centering\scriptsize\setlength\tabcolsep{2pt}\renewcommand{\arraystretch}{0.78}
\caption{Representation intervention and control studies (an effect exceeding the reported uncertainty establishes causal sensitivity; direction-selective sensitivity additionally requires a matched control direction, while matched features or tokens establish component selectivity only; a concept-specific causal contribution requires a mediation or causal-abstraction design that no reviewed medical study reports, \S\ref{sec:causal-claims}; column Ctrl records whether any such control was reported, not which kind) (table 1 of 3). Core medical studies only; columns and symbols as in Table~S1. Ctrl: matched control reported for the intervention (\cmark\ yes, \xmark\ no, n/a = editing, unlearning, erasure, batch-correction or embedding-shift design for which a matched-direction control is not the relevant comparison).}
\label{tab:intervention}
\begin{tabular}{@{}L{2.4cm}L{1.8cm}L{1.1cm}L{0.8cm}L{0.8cm}C{0.5cm}C{0.6cm}C{0.6cm}C{0.6cm}C{0.6cm}L{5.0cm}@{}}
\toprule
Study & Model(s) & Modality & Locus & Prov. & Ev. & Base. & Multi. & 2nd & Ctrl & Main finding \\
\midrule
\multicolumn{11}{@{}l}{\emph{Ablation and knockout}} \\
\cite{chen2026vihd} (2026) \newline {\tiny\texttt{chen2026vihd}} & Hulu-Med-4B; Lingshu-7B & multi & Dec & none & I & \xmark & \cmark & tra & \cmark & Masking the 10\% highest-attention visual tokens inside the visually dominant decoder layers and scoring calibrated semantic entropy detects hallucinations at AUC 71.6-78.9 vs 68.9-73.0 for the best introspective baseline (two medical MLLMs, three VQA benchmarks); random or low-attention masking produced little change in the reported point estimates, single runs (+7.41 AUC for targeted masking on VQA-RAD) \\
\cite{nooralahzadeh2026sparse} (2026) \newline {\tiny\texttt{nooralahzadeh2026sparse}} & frozen 3D CT encoders & CT & Enc & label & I & \xmark & \pmark & dec & \xmark & Zeroing a finding's ~10 channels collapses its own AUROC while off-target findings stay stable (specificity control) \\
\multicolumn{11}{@{}l}{\emph{Steering vectors and representation engineering}} \\
\cite{bouzid2025insights} (2025) \newline {\tiny\texttt{bouzid2025insights}} & MAIRA-2 & CXR & Dec & none & I & \xmark & \xmark & dis & \xmark & Matryoshka-SAE features; steering with mixed success \\
\cite{farina2026translating} (2026) \newline {\tiny\texttt{farina2026translating}} & three CXR VLMs & CXR & Dec & label & I & \xmark & \pmark & -- & \cmark & CAA steering with random/reverse controls; gains in one of three models \\
\cite{nooralahzadeh2026universal} (2026) \newline {\tiny\texttt{nooralahzadeh2026universal}} & RadVLM; LLaVA-Rad; CheXOne & CXR & Dec & none & I & \xmark & \cmark & dis & \xmark & SAE steering improves clinical composite; boost directions shared; suppressors model-specific; no random-direction null \\
\cite{restrepo2026cone} (2026) \newline {\tiny\texttt{restrepo2026cone}} & BiomedCLIP; MedSigLIP & CXR; derm; retina & Joint & none & I & \cmark & \cmark & geo & n/a & Post-hoc shift along the gap vector on frozen encoders; moderate closure helps; full collapse hurts \\
\cite{zeng2026easylens} (2026) \newline {\tiny\texttt{zeng2026easylens}} & MedGemma 1.5; MedGemma; Lingshu; LLaVA-Med; RadFM & CT; endoscopy; CXR & Enc & label & I & \xmark & \cmark & -- & \xmark & Training-free amplification of lesion-relevant patch tokens (prototypes from lesion masks) before the LLM raises subtle-lesion scores; five backbones on ReX, MedGemma 1.5 on LIDC/Abdomen \\
\bottomrule
\end{tabular}
\end{table*}

\begin{table*}[!tbp]
\centering\scriptsize\setlength\tabcolsep{2pt}\renewcommand{\arraystretch}{0.78}
\caption{Representation intervention and control studies (an effect exceeding the reported uncertainty establishes causal sensitivity; direction-selective sensitivity additionally requires a matched control direction, while matched features or tokens establish component selectivity only; a concept-specific causal contribution requires a mediation or causal-abstraction design that no reviewed medical study reports, \S\ref{sec:causal-claims}; column Ctrl records whether any such control was reported, not which kind) (table 2 of 3; continued from Table~\ref{tab:intervention}). Columns and symbols as in Table~S1.}
\label{tab:intervention2}
\begin{tabular}{@{}L{2.4cm}L{1.8cm}L{1.1cm}L{0.8cm}L{0.8cm}C{0.5cm}C{0.6cm}C{0.6cm}C{0.6cm}C{0.6cm}L{5.0cm}@{}}
\toprule
Study & Model(s) & Modality & Locus & Prov. & Ev. & Base. & Multi. & 2nd & Ctrl & Main finding \\
\midrule
\multicolumn{11}{@{}l}{\emph{Activation patching and causal tracing}} \\
\cite{sadanandan2026mechanistically} (2026) \newline {\tiny\texttt{sadanandan2026mechanistically}} & MedGemma & CXR & Dec & none & I & \xmark & \xmark & dis & \cmark & Patching a single SAE feature recovers more of the decision margin than random control features; guided LoRA improves paraphrase consistency; early layers beat dictionary-selected middle layers \\
\multicolumn{11}{@{}l}{\emph{Post-hoc transformation, editing and unlearning}} \\
\cite{weber2023posthoc} (2023) \newline {\tiny\texttt{weber2023posthoc}} & CXR Foundation; CheSS & CXR & Enc & label & I & \pmark & \cmark & dec & n/a & QR-orthogonalisation of frozen CXR embeddings against age, sex and race sets every protected-feature coefficient on the pathology logit to zero and attribute recovery to chance (CFM sex AUC 0.979 to 0.507) at a macro-AUC cost of -1.39\% (CFM/MIMIC) to +1.83\% (CheSS/CheXpert); subgroup AUC gaps narrow \\
\cite{murchan2024deep} (2024) \newline {\tiny\texttt{murchan2024deep}} & CTransPath; RetCCL (vs ImageNet ResNet50) & pathology & Enc & label & I & \pmark & \cmark & dec & n/a & ComBat on frozen patch embeddings (site as batch) drops A-MIL site prediction from AUROC 0.955-0.977 to 0.506-0.591 in 5 of 6 model-cohort pairs and race from 0.835 to 0.471, while MSI (0.667->0.669, p=0.95) holds \\
\cite{sadashivaiah2024explaining} (2024) \newline {\tiny\texttt{sadashivaiah2024explaining}} & CXR VLM joint space & CXR & Joint & text & I & \xmark & \xmark & -- & \xmark & Conceptual counterfactuals: embedding moved along report-derived textual concept directions in a CXR VLM joint space to change the prediction \\
\cite{carloni2025pathology} (2025) \newline {\tiny\texttt{carloni2025pathology}} & Phikon; UNI; Virchow; Prov-GigaPath; H-Optimus-0 & pathology & Enc & label & I & \xmark & \pmark & geo & n/a & ScanGen contrastive projection of frozen embeddings (same-specimen attraction across scanners) lowers prediction CoV across scanners by 8-69\% over five FMs and keeps or improves EGFR AUC \\
\cite{folco2025semantic} (2025) \newline {\tiny\texttt{folco2025semantic}} & RAD-DINO (vs DINOv2, ViT-S); BioLORD, Meditron & CXR & Joint & text & I & \cmark & \cmark & geo & n/a & Transporting frozen image embeddings into the CLIP text space with the same map gives zero-shot classification without multimodal training, beating CLIP on six of seven CXR tasks and approaching MedCLIP; general encoders beat RAD-DINO on most tasks \\
\cite{zhang2025hospital} (2025) \newline {\tiny\texttt{zhang2025hospital}} & UNI; UNI2-H; Prov-GigaPath; CONCH; TITAN; MUSK; Phikon-v2 & pathology & Enc & label & I & \pmark & \pmark & dec & n/a & A gradient-reversal projection head on frozen embeddings drops hospital-source AUC from 0.95/0.96/0.91 to 0.58/0.51/0.49 for UNI, UNI2-H and Virchow while IDC-vs-ILC accuracy holds (0.93/0.92/0.91); t-SNE clusters dissolve \\
\bottomrule
\end{tabular}
\end{table*}

\begin{table*}[!tbp]
\centering\scriptsize\setlength\tabcolsep{2pt}\renewcommand{\arraystretch}{0.78}
\caption{Representation intervention and control studies (an effect exceeding the reported uncertainty establishes causal sensitivity; direction-selective sensitivity additionally requires a matched control direction, while matched features or tokens establish component selectivity only; a concept-specific causal contribution requires a mediation or causal-abstraction design that no reviewed medical study reports, \S\ref{sec:causal-claims}; column Ctrl records whether any such control was reported, not which kind) (table 3 of 3; continued from Table~\ref{tab:intervention}). Columns and symbols as in Table~S1.}
\label{tab:intervention3}
\begin{tabular}{@{}L{2.4cm}L{1.8cm}L{1.1cm}L{0.8cm}L{0.8cm}C{0.5cm}C{0.6cm}C{0.6cm}C{0.6cm}C{0.6cm}L{5.0cm}@{}}
\toprule
Study & Model(s) & Modality & Locus & Prov. & Ev. & Base. & Multi. & 2nd & Ctrl & Main finding \\
\midrule
\multicolumn{11}{@{}l}{\emph{Post-hoc transformation, editing and unlearning}} \\
\cite{arora2026semantic} (2026) \newline {\tiny\texttt{arora2026semantic}} & CPath-CLIP; H-optimus-0 & pathology & Joint & text & I & \xmark & \cmark & geo & n/a & Scoring frozen CPath-CLIP patch embeddings against histological text anchors instead of image prototypes lifts human-to-canine tumour AUC from 63.96 to 78.28, while canine-specific anchors give 64.79; an off-target PatchCamelyon control is reported \\
\cite{balogh2026regularization} (2026) \newline {\tiny\texttt{balogh2026regularization}} & BiomedCLIP; PubMedCLIP & CXR & Joint & text & I & \xmark & \pmark & geo & n/a & A similarity-level regulariser on the frozen encoders' text projection head cuts the vulnerability index 34-84\% and raises detection AUC to 0.729 negation and 0.687 laterality, but held-out directional gain is only +0.005 \\
\cite{donle2026featmap} (2026) \newline {\tiny\texttt{donle2026featmap}} & pathology and brain-MRI FMs & pathology; MRI & Enc & label & I & \xmark & \cmark & geo & n/a & Acquisition signature modelled as affine manifold distortion; one correction per FM harmonises embeddings \\
\cite{garciahidalgo2026embedding} (2026) \newline {\tiny\texttt{garciahidalgo2026embedding}} & eight encoders incl. BiomedCLIP; PLIP; UNI; Virchow & multi & Enc & none & I & \cmark & \cmark & geo & n/a & Post-hoc whitening degrades downstream performance; gentle spectral standardisation gives a marginal gain (about +0.005 AUC) \\
\cite{kim2026encoding} (2026) \newline {\tiny\texttt{kim2026encoding}} & six CXR models (nine with three CXR-specialised additions) & CXR & Enc & label & I & \cmark & \cmark & dec & n/a & Residualising the attribute out changes the AUROC of refitted finding probes by <0.001: linearly redundant for an optimal refitted readout (fixed-head use untested) \\
\cite{kmen2026robust} (2026) \newline {\tiny\texttt{kmen2026robust}} & 20 pathology FMs (PathoROB; incl. UNI2-h; Virchow2; CONCH; H-optimus-0; Prov-GigaPath; Phikon) & pathology & Enc & label & I & \xmark & \cmark & dec & n/a & ComBat removal of centre signature from embeddings raises robustness index +28\% yet does not improve downstream generalisation; centre and biology share linear directions \\
\cite{li2026acquisition} (2026) \newline {\tiny\texttt{li2026acquisition}} & RETFound (vs DINOv3 ViT-S/B, ConvNeXt; DINOv2 ViT-S/L) & retina & Enc & label & I & \cmark & \pmark & -- & n/a & Training-light debiasing of the same frozen embeddings (per-camera harmonisation, INLP): device decodability falls only from 0.95 to 0.70-0.86, the ~3-year age-gap device shift persists, and INLP k=20 costs age MAE 7.6 to 14 y \\
\bottomrule
\end{tabular}
\end{table*}
